\honest{Typicality and Asymmetrical Similarity}{Eliezer Yudkowsky}{2008}

Is a robin or an ostrich the more typical bird? A desk chair, a rocking chair or a beanbag the more typical chair? Most people say the robin and the desk chair, and psychologists study this as ``typicality effects.'' People confirm faster that a robin is a bird than that a penguin is, and such reaction times agree with people's direct ratings, from 1 to 10, of how well an example fits a category. I am ``reasonably sure'' my source for this is Lakoff, though my books are still in boxes.

Typicality may be a heuristic, so I look for a bias it produces. ``98 is approximately 100'' sounds more natural than ``100 is approximately 98,'' and I cite Sadock (1977). \nb{Sadock's paper does not contain this comparison. The experiment behind it is Rosch's work on reference points: people put the more typical item in the reference slot of phrases like ``X is virtually Y.''} People also rate Mexico as more similar to the United States than the reverse. \nb{In Tversky and Gati's table the ratings are 7.65 and 6.46 on a scale of 1 to 20. Across their 21 pairs of countries the average difference was 0.42, and six pairs went the other way.}

In Rips's study, people expected a disease to spread from robins to ducks more readily than from ducks to robins. Any difference that would slow a disease one way would ``in a pragmatic sense'' slow it the other way too. \nb{A Bayesian model by Heit disagrees: if atypical species are believed to have many features of their own, a new property of ducks is more likely to be theirs alone.} One could rationalize, say, that robins have more neighboring species to spread a disease to, but do not try too hard to rationalize answers from subjects ``who didn't even realize there was a comparison going on,'' and remember Mexico and 98. The simpler explanation is that people use similarity as a stand-in for probability, and similarity is asymmetric. \nb{Rips measured similarity symmetrically and found a separate effect of how typical the species in the premise was. On Rips's analysis the asymmetry comes from typicality, not from asymmetric similarity.}

Kansas lies near the center of the United States, so it is probably closer than Alaska to most places. But that does not make Kansas closer to Alaska than Alaska is to Kansas; people reason as if closeness were a property of Kansas and distance a property of Alaska. \nb{A 1980 study by Sadalla and colleagues found the same asymmetry in judged distances between reference points and other places.}

So Aristotle's categories, with necessary and jointly sufficient properties, are a poor model of how people think (has science's view changed in 2,350 years? who would have thought), and ``Statements of set membership can be more or less true,'' which is not the same as more or less probable. \nb{Typicality does not show this. People rate 3 as a more typical odd number than 4,284,647, though both are plainly odd.} One more reason not to pretend that anyone treats words as Aristotelian classes.
